\chapter{Обзор предметной области и постановка задачи}

\section{Проблема траекторного управления квадрокоптером}

В работе рассматривается задача движения квадрокоптера вдоль заранее заданной пространственной траектории. На практике здесь важна не только сама сходимость к маршруту, но и качество движения после выхода на него: поперечная ошибка должна оставаться малой, продольная синхронизация не должна разрушаться при изменении режима, а скорость продвижения вдоль кривой должна быть достаточно высокой, но безопасной для текущей геометрии траектории.

Классические постановки такой задачи обычно опираются на аналитический регулятор слежения, в котором заранее задается параметрическая скорость движения вдоль кривой. Такой подход хорошо изучен и дает устойчивый базовый контур управления. В частности, близкая логика используется в работах по нелинейному и адаптивному управлению сложными динамическими системами, а также в исследованиях по траекторному управлению квадрокоптерами на основе геометрических и выходных представлений~\footcite{miroshnik2000,borisov2021,kim2023,kim2024article}.

Проблема состоит в том, что фиксированная скорость $V^{\ast}$ почти всегда выбирается с запасом. На прямолинейных участках и в спокойных режимах объект может двигаться быстрее без заметного роста ошибки, а на участках с высокой кривизной или при неудачном начальном возмущении та же самая скорость уже приводит к ухудшению качества слежения. Из-за этого возникает практическая задача адаптивного выбора $V^{\ast}$ по текущему состоянию системы.

\section{Базовый аналитический контур}

В качестве нижнего уровня управления в проекте используется алгоритм согласованного траекторного управления по выходу, основанный на работах С.~А.~Кима и соавторов~\footcite{kim_algorithms,kim2024article}. Его ключевая идея состоит в том, что управление строится не напрямую по декартовым координатам, а по ошибкам относительно ближайшей точки опорной траектории. Такой переход делает связь между динамикой квадрокоптера и геометрией маршрута более прозрачной и позволяет отдельно анализировать продольную и поперечную составляющие ошибки.

С инженерной точки зрения этот подход удобен по трем причинам. Во-первых, он дает устойчивый базовый режим, который можно исследовать независимо от интеллектуального модуля. Во-вторых, он хорошо переносится в вычислительную реализацию и допускает автоматизированную проверку на типовых траекториях. В-третьих, именно поверх такого контура можно безопасно подключать модуль выбора скорости, не передавая нейросети управление всем объектом целиком.

\section{Обзор исследований по теме}

\subsection{Аналитические методы траекторного управления БПЛА}

Теоретическую основу нелинейного управления сложными динамическими системами составляет монография Мирошника, Никифорова и Фрадкова~\footcite{miroshnik2000}, в которой систематически изложены методы управления по выходу, адаптивного управления и синтеза нелинейных регуляторов. На эти результаты опирается математический аппарат базового контура данной работы.

Применительно к квадрокоптеру траекторное управление по выходу исследовалось в серии работ Борисова, Кима, Пыркина и соавторов. В статье 2021 года предложен подход к робастному траекторному управлению на основе геометрического представления ошибок~\footcite{borisov2021}. В работе 2023 года рассматривается режим динамического позиционирования~\footcite{kim2023}. Алгоритм согласованного траекторного управления~\footcite{kim2024article}, лежащий в основе данной ВКР, опубликован в 2024 году и описывает метод, при котором продольная и поперечная ошибки анализируются в сопутствующей системе координат и связываются единым законом управления.

\subsection{Интеграция нейронных сетей в системы управления БПЛА}

Современные работы по управлению БПЛА показывают, что интеграция методов машинного обучения обычно идёт по одному из двух путей.

Первый~--- нейросеть заменяет часть классического регулятора или напрямую вырабатывает управляющее воздействие. В~\cite{kownacki2024} нейронные сети применяются как прямые контроллеры в задаче позиционного и траекторного управления holonomic UAV. Такой подход даёт высокую гибкость, но требует особенно аккуратной верификации устойчивости и переносимости.

Второй путь ближе к задаче данной ВКР: методы машинного обучения используются не для полной замены аналитического контура, а для улучшения отдельных компонентов. Работа~\cite{madruga2021} показывает, как нейронная сеть может компенсировать аэродинамические эффекты, разгружая классический контроллер от моделирования трудно поддающихся измерению возмущений. В~\cite{shi2025} глубокое обучение с подкреплением совмещается с нелинейным MPC для планирования и сопровождения траектории, что демонстрирует жизнеспособность гибридных подходов. Применение нейросетей для задач планирования маршрута рассмотрено в~\cite{akshya2024}.

\subsection{Оптимизация скорости движения БПЛА}

Отдельный интерес представляют исследования, в которых скорость движения рассматривается как самостоятельный объект оптимизации. В~\cite{wu2023} нейросетевой модуль совмещается с методом роя частиц (PSO) для управления скоростью БПЛА при посадке на подвижную платформу. В~\cite{wang2025} предложена иерархическая схема, в которой верхний уровень занимается оптимизацией скорости, а нижний~--- сопровождением трассы в условиях присутствия других агентов. Общий вывод этих работ согласуется с логикой данного проекта: допустимая скорость зависит не только от команды оператора, но и от локальной геометрии траектории, текущей ошибки и динамического состояния системы, и именно поэтому задача её выбора нуждается в адаптивном решении.

\subsection{Алгоритмы обучения с подкреплением, используемые в работе}

Для построения трёх альтернативных моделей выбора скорости в данной работе использованы три алгоритма, каждый из которых имеет чёткое теоретическое обоснование. SAC (Soft Actor-Critic)~\footcite{haarnoja2018sac} вводит энтропийный регуляризатор, предотвращающий преждевременную детерминизацию политики. TD3 (Twin Delayed DDPG)~\footcite{fujimoto2018td3} использует пару критиков и задержку обновления актора для борьбы с переоценкой Q-функции. PPO~\footcite{schulman2017ppo} ограничивает шаг изменения политики коэффициентом клиппинга, что делает обучение стабильнее при разнородных данных. Все три алгоритма применяются в данной работе в режиме offline-обучения по готовому датасету, а не в классическом режиме взаимодействия со средой.

\section{Постановка задачи ВКР}

С учетом рассмотренного обзора задача выпускной квалификационной работы формулируется следующим образом: необходимо разработать и исследовать программную систему, в которой классический контур траекторного управления квадрокоптером дополняется модулем интеллектуального выбора параметрической скорости движения вдоль траектории.

Для достижения этой цели решаются следующие подзадачи:
\begin{enumerate}
    \item реализовать и проверить базовый аналитический регулятор траекторного слежения;
    \item сформировать датасет состояний и соответствующих им допустимых скоростей движения;
    \item обучить несколько моделей машинного обучения для предсказания значения $V^{\ast}$;
    \item встроить обученные модели в контур моделирования с ограничениями, гарантирующими безопасную работу;
    \item сравнить модели между собой и с базовым режимом на наборе воспроизводимых тестовых траекторий.
\end{enumerate}

В работе сознательно выбран умеренный по риску подход: нейросетевой модуль не заменяет аналитический регулятор, а выступает как надстройка над ним. Это решение продиктовано не только удобством реализации, но и инженерной логикой. Если нижний уровень уже обеспечивает устойчивое слежение, то верхнему уровню достаточно решать более локальную и практически полезную задачу: когда можно ускориться, а когда лучше снизить темп, чтобы не потерять качество траектории.

\section{Используемая литература и ее роль в работе}

Использованную литературу можно условно разделить на три группы. Первая группа включает теоретические источники по нелинейному и согласованному управлению~\footcite{miroshnik2000,kim_algorithms,kim2024article}. На них опирается математическая постановка базового контура и описание регулируемых переменных.

Вторая группа объединяет прикладные публикации по управлению квадрокоптерами и БПЛА~\footcite{borisov2021,kim2023,madruga2021,kownacki2024}. Эти работы важны для понимания того, какие архитектурные решения действительно применимы к динамическим летательным объектам, а какие остаются слишком общими.

Третья группа относится к методам машинного обучения и обучения с подкреплением, используемым при построении модуля выбора скорости~\footcite{kingma2015adam,schulman2017ppo,haarnoja2018sac,fujimoto2018td3,shi2025}. Она задает основу для выбора архитектур моделей, алгоритмов обучения и критериев сравнения.

\section{Выводы по главе}

В первой главе показано, что задача адаптивного выбора скорости движения естественно возникает внутри классической схемы траекторного управления квадрокоптером. Аналитический регулятор обеспечивает устойчивую основу, а методы машинного обучения позволяют повысить эффективность движения без полной замены базового контура. Такой подход определяет структуру всей дальнейшей работы: во второй главе рассматриваются модели, методы и программная реализация, а в третьей главе приводятся результаты вычислительных экспериментов.
